% chktex-file 19
% chktex-file 35
% chktex-file 36
 \documentclass[12pt]{article}
 \usepackage[left=2cm,right=2cm,top=2cm,bottom=2cm,bindingoffset=0cm]{geometry}
 \usepackage{amssymb}
 \usepackage{amsmath}
 \usepackage{amsthm}
 \usepackage[T1]{fontenc}
 \usepackage[utf8]{inputenc}
 \usepackage[colorlinks=true, urlcolor=blue, linkcolor=blue, citecolor=blue,
     pdfborder={0 0 0}]{hyperref}
\usepackage{enumitem}

 \newtheorem{theorem}{Theorem}[section]
 \newtheorem{lemma}[theorem]{Lemma}
 \newtheorem{proposition}[theorem]{Proposition}
 \newtheorem{corollary}[theorem]{Corollary}
 \theoremstyle{definition}
 \newtheorem{definition}[theorem]{Definition}
 \newtheorem{example}[theorem]{Example}
 \newtheorem{exercise}[theorem]{Exercise}
 \theoremstyle{remark}
 \newtheorem{remark}[theorem]{Remark}

 \newcommand{\lo}[1]{\ell_1^0(#1)}
 \newcommand{\ellone}[1]{\ell_1(#1)}
 \newcommand{\norm}[1]{\lVert #1\rVert}
 \newcommand{\supp}{\operatorname{supp}}
 \newcommand{\Lin}{\operatorname{span}}
 \newcommand{\Lim}{\operatorname{Lim}}
 \newcommand{\NN}{\mathbb{N}}

\begin{document}

\begin{center}
    \Large \textbf{Set theory for the paper ``Extremely projective normed spaces''}\\[0.3cm]
    \large A crash course\\[0.5cm]
\end{center}

\thispagestyle{empty}

\medskip
\noindent\textbf{What this is.} The paper ``Extremely projective normed
spaces'' uses, in its sections 6 and 7, a handful of tools from set theory:
ordinals and transfinite recursion, the first uncountable ordinal $\omega_1$,
closed unbounded and stationary subsets of $\omega_1$, and one theorem of
Solovay. These notes explain exactly those tools, assuming only naive set
theory and the notion of cardinality. Section~\ref{SecOrdinals} introduces
well-orders, ordinals and transfinite recursion; section~\ref{SecCardinals}
reviews cardinals and the axiom of choice; section~\ref{SecOmegaOne} treats
$\omega_1$ and its closed unbounded sets; section~\ref{SecReading} walks
through the set-theoretic steps of the paper one by one; and
section~\ref{SecGlossary} is a glossary. Exercises come with hints; the ones
marked with an asterisk are used in the paper and are worth doing in full.
For a textbook treatment see Jech, \emph{Set theory}, chapters 2, 3, 5 and 8,
or, for a gentler start, Hrbacek and Jech, \emph{Introduction to set theory}.
\medskip

%-------------------------------------------------------------------------------
\section{Well-orders, ordinals and transfinite recursion}\label{SecOrdinals}

\subsection{Well-orders}

\begin{definition} A \emph{well-order} is a linearly ordered set $(W,<)$ in
    which every non-empty subset has a least element.
\end{definition}

\begin{example} $\NN$ with its usual order is a well-order. So is any finite
    linear order. So is $\NN\cup\{\infty\}$ with $\infty$ placed after all
    natural numbers: a non-empty subset either contains a natural number, and
    then has a least one, or is $\{\infty\}$. So is ``two copies of $\NN$ one
    after the other'', the set $\{0,1\}\times\NN$ ordered lexicographically
    with the first coordinate dominant. The sets $\mathbb{Z}$, $\mathbb{Q}$ and
    $[0,1]$ with their usual orders are not well-orders: $\mathbb{Z}$ has no
    least element, and $(0,1]\subset[0,1]$ has none.
\end{example}

Two consequences of the definition are used all the time. First, a well-order
has no infinite strictly decreasing sequence $x_0>x_1>x_2>\cdots$, since the
set $\{x_n:n\in\NN\}$ would have no least element. Second, the
\emph{principle of induction}: if a property $\Phi(x)$ fails for some
$x\in W$, then it fails for a least $x$, so to prove $\Phi$ for all $x$ it
suffices to prove $\Phi(x)$ under the assumption that $\Phi(y)$ holds for all
$y<x$.

For $x\in W$ the \emph{initial segment} $W_{<x}=\{y\in W:y<x\}$ is again a
well-order. The basic structure theorem, which we state without proof, is
that well-orders are linearly ordered by ``being an initial segment'':

\begin{theorem}\label{ComparabilityOfWellOrders} Let $V$ and $W$ be
    well-orders. Then exactly one of the following holds: $V$ and $W$ are
    order isomorphic; $V$ is order isomorphic to $W_{<x}$ for a unique
    $x\in W$; $W$ is order isomorphic to $V_{<y}$ for a unique $y\in V$. In
    each case the isomorphism is unique.
\end{theorem}

\subsection{Ordinals}

Ordinals are canonical representatives of well-orders, chosen so that the
order relation becomes membership. Following von Neumann, put
\begin{gather*}
0=\varnothing,\quad 1=\{0\},\quad 2=\{0,1\},\quad\ldots,\quad
n=\{0,1,\ldots,n-1\},\quad\ldots,\\
\omega=\{0,1,2,\ldots\}=\NN,\quad \omega+1=\omega\cup\{\omega\},
\end{gather*}
and so on: each ordinal is the set of all smaller ordinals, ordered by
membership. Formally, an \emph{ordinal} is a set $\alpha$ which is well-ordered
by $\in$ and \emph{transitive}, meaning that every element of an element of
$\alpha$ is an element of $\alpha$. We write $\beta<\alpha$ for
$\beta\in\alpha$. The facts to remember are the following.
\begin{enumerate}[label = (\alph*)]
    \item Every well-ordered set is order isomorphic to exactly one ordinal,
          called its \emph{order type}. Ordinals are therefore ``the''
          well-orders, and theorem~\ref{ComparabilityOfWellOrders} says that
          any two ordinals $\alpha,\beta$ satisfy exactly one of
          $\alpha<\beta$, $\alpha=\beta$, $\beta<\alpha$, where $\alpha<\beta$
          also means $\alpha\subsetneq\beta$.
    \item Every element of an ordinal is an ordinal, and the elements of
          $\alpha$ are exactly the ordinals $\beta<\alpha$. Thus,
          $\alpha=\{\beta:\beta<\alpha\}$. This is why the paper writes
          $\ellone{\delta}$ for the families supported in
          $\{\beta:\beta<\delta\}$: the ordinal $\delta$ \emph{is} that set.
    \item Every set of ordinals has a least element (it is a subset of a
          well-order), and a supremum, namely its union, which is again an
          ordinal. There is no set of all ordinals: for every set $M$ of
          ordinals, $\sup M+1$ is an ordinal not in $M$.
    \item The \emph{successor} of $\alpha$ is $\alpha+1=\alpha\cup\{\alpha\}$,
          the least ordinal above $\alpha$. An ordinal is a \emph{successor
          ordinal} if it is of the form $\alpha+1$, and a \emph{limit
          ordinal} if it is not a successor and not $0$. (Some authors count
          $0$ as a limit ordinal; the paper explicitly does not.) A limit
          ordinal $\delta$ is the supremum, that is, the union, of the ordinals
          below it: $\delta=\bigcup_{\beta<\delta}\beta$. Every ordinal is
          $0$, a successor or a limit.
\end{enumerate}

\begin{example}\label{SmallOrdinals} The finite ordinals are the natural
    numbers. The first infinite ordinal is $\omega$, the order type of $\NN$,
    a limit. Then come $\omega+1,\omega+2,\ldots$, successors; their supremum
    $\omega+\omega=\omega\cdot 2$, the order type of two copies of $\NN$ one
    after the other, is the second limit ordinal; then
    $\omega\cdot 2+1,\ldots,\omega\cdot 3,\ldots,\omega\cdot\omega=\omega^2$, the
    order type of $\NN\times\NN$ ordered lexicographically; and so on. All of
    these are countable sets. The limit ordinals among them are exactly the
    non-zero multiples of $\omega$.
\end{example}

\begin{exercise} Show that $\omega$ and $\omega+1$ are not order isomorphic.
    (Hint: one of them has a largest element.) Show that $\omega\cdot 2$ is
    not order isomorphic to $\omega$. (Hint: in $\omega\cdot 2$ the element
    $\omega$ has infinitely many predecessors but no immediate predecessor.)
\end{exercise}

\begin{exercise} Let $(W,<)$ be a linear order with no infinite strictly
    decreasing sequence. Show that $W$ is a well-order. (Hint: if a non-empty
    subset $M$ has no least element, choose $x_0\in M$, then $x_1<x_0$ in
    $M$, and so on; this uses countable choice, see
    section~\ref{SecChoice}.)
\end{exercise}

\subsection{Transfinite induction and recursion}

The principle of induction for well-orders becomes, for ordinals, the
principle of \emph{transfinite induction}: to prove $\Phi(\alpha)$ for all
ordinals $\alpha$ (or for all $\alpha\leq\lambda$) it suffices to prove
$\Phi(\alpha)$ assuming $\Phi(\beta)$ for all $\beta<\alpha$. In practice one
splits into three cases: $\alpha=0$; $\alpha=\beta+1$, where one assumes
$\Phi(\beta)$; and $\alpha$ a limit, where one assumes $\Phi(\beta)$ for all
$\beta<\alpha$.

More important for the paper is \emph{transfinite recursion}: an object
$F(\alpha)$ may be defined for all ordinals $\alpha<\lambda$ by a rule which
produces $F(\alpha)$ from the family ${(F(\beta))}_{\beta<\alpha}$. The
theorem of transfinite recursion says that such a rule determines a unique
function $F$ on $\lambda$. Again one usually gives the rule in three parts:
what $F(0)$ is, how $F(\beta+1)$ is obtained from $F(\beta)$, and how
$F(\alpha)$ is obtained from ${(F(\beta))}_{\beta<\alpha}$ at a limit
$\alpha$, typically as a union or a limit.

\begin{example}[the chains of closed pairs in the paper] In the proof of
    theorem~6.7 of the paper a family ${((A_\alpha,B_\alpha))}_{\alpha\leq\kappa}$
    is defined by transfinite recursion: $(A_0,B_0)=(\varnothing,\varnothing)$;
    $(A_{\alpha+1},B_{\alpha+1})$ is obtained from $(A_\alpha,B_\alpha)$ by
    adding one point and closing off; at a limit $\alpha$ the pair
    $(A_\alpha,B_\alpha)$ is the union of the previous pairs. In the proof of
    theorem~6.1 the operators $V_\delta$ are defined by recursion on
    $\delta<\lambda$: the rule producing $V_\delta$ uses $W_\delta$ and
    $T_\delta$, which depend only on the $V_{\delta'}$ with $\delta'<\delta$.
    The accompanying claim $T_\delta(P_\delta)=\lo{\Gamma_\delta}$ is then
    proved by transfinite induction on $\delta$, with the limit case handled
    by the continuity hypothesis.
\end{example}

\subsection{Filtrations}

The paper calls a family ${(\Gamma_\delta)}_{\delta\leq\lambda}$ of subsets
of a set $\Gamma$ a \emph{filtration} if it is increasing,
$\Gamma_0=\varnothing$, $\Gamma_\lambda=\Gamma$, and it is \emph{continuous}:
$\Gamma_\delta=\bigcup_{\delta'<\delta}\Gamma_{\delta'}$ for every limit
ordinal $\delta\leq\lambda$. Thus, a filtration grows only at successor
steps, by the \emph{blocks} $N_\delta=\Gamma_{\delta+1}\setminus\Gamma_\delta$,
and at limit steps it merely collects what has been built. Continuity is a
genuine condition: the family given by $\Gamma_n=\{0,\ldots,n-1\}$ for
$n<\omega$, $\Gamma_\omega=\NN\cup\{\star\}$ and $\Gamma_{\omega+1}=\NN\cup\{\star,\star\star\}$
is increasing but not continuous at $\omega$, because the point $\star$
appears at the limit stage without ever appearing in a block.

\begin{exercise}[*] Prove by transfinite induction on $\delta\leq\lambda$ that
    for a filtration $\Gamma_\delta=\bigsqcup_{\delta'<\delta}N_{\delta'}$,
    a disjoint union. (Hint: the successor step is the definition of
    $N_\delta$; the limit step is continuity.) Deduce that every $x\in\ellone{\Gamma}$
    satisfies $\norm{x}=\sum_{\delta'<\lambda}\norm{x|_{N_{\delta'}}}$, and
    that only countably many of the summands are non-zero.
\end{exercise}

\begin{exercise}[*] Let $C\subset\omega_1$ be a set and
    ${(c_\alpha)}_{\alpha<\mu}$ its increasing enumeration (the elements of
    $C$ listed in increasing order; see section~\ref{SecCardinals} for why
    the index set is an ordinal). Show that $\Gamma_\alpha=c_\alpha$ defines a
    continuous family if and only if $C$ is closed under suprema of its
    bounded subsets. This is how the paper turns a closed unbounded set into
    a filtration in the proof of theorem~7.2.
\end{exercise}

%-------------------------------------------------------------------------------
\section{Cardinals and the axiom of choice}\label{SecCardinals}

\subsection{Cardinals as initial ordinals}

Two sets have the same cardinality if there is a bijection between them. To
have a canonical representative of each cardinality one uses ordinals again:
the \emph{cardinal} $|X|$ of a set $X$ is the least ordinal which is in
bijection with $X$. For this to make sense $X$ must be in bijection with some
ordinal at all, that is, $X$ must admit a well-order; that every set does is
the \emph{well-ordering theorem}, equivalent to the axiom of choice, see
section~\ref{SecChoice}. An ordinal which is the cardinal of some set, that
is, which is not in bijection with any smaller ordinal, is called an
\emph{initial ordinal}. The finite ordinals are initial, and $\omega$ is
initial; it is also written $\aleph_0$ when regarded as a cardinal. The
ordinals $\omega+1$, $\omega\cdot 2$, $\omega^2$ of example~\ref{SmallOrdinals}
are not initial: they are countable, hence in bijection with $\omega$.

\begin{definition} $\omega_1$ is the least uncountable ordinal, that is, the
    least ordinal not in bijection with $\omega$. Regarded as a cardinal it
    is written $\aleph_1$. Likewise $\omega_2=\aleph_2$ is the least ordinal
    of cardinality greater than $\aleph_1$, and so on.
\end{definition}

Such ordinals exist: the order types of well-orderings of subsets of $\NN$
form a set of ordinals, all countable, and every countable ordinal is among
them; the supremum of this set plus one is uncountable, and the least
uncountable ordinal exists because every set of ordinals has a least element.
By definition, the elements of $\omega_1$ are exactly the countable ordinals:
\[
\omega_1=\{\alpha:\alpha\text{ is a countable ordinal}\}.
\]
This double role is the source of most confusion for beginners. The symbol
$\omega_1$ denotes both a well-ordered set, the set of countable ordinals in
their natural order, and a cardinality $\aleph_1$, the cardinality of that
set. When the paper says ``let $\Gamma=\omega_1$'' it uses the first role: the
index set is well-ordered, and the subspaces $\ellone{\delta}$,
$\delta<\omega_1$, are its initial segments. When it says
``$|\Gamma|=\aleph_1$'' it uses the second role.

The \emph{continuum hypothesis}, the statement $2^{\aleph_0}=\aleph_1$, is
independent of the usual axioms and is \emph{not} used anywhere in the paper.
Nothing in the paper depends on how large $2^{\aleph_0}$ or $2^{\aleph_1}$
is; the latter only appears as the number of subsets of $\omega_1$.

\subsection{Cardinal arithmetic used in the paper}

Only the most basic facts are needed, all of which hold for an infinite
cardinal $\kappa$:
\begin{enumerate}[label = (\alph*)]
    \item $\kappa\cdot\aleph_0=\kappa$ and $\kappa+\kappa=\kappa$; more
          generally $\kappa\cdot\mu=\max(\kappa,\mu)$ for infinite $\kappa,\mu$;
    \item the set of finite subsets of a set of cardinality $\kappa$ has
          cardinality $\kappa$, and so has the set of countable subsets of a
          countable set... which is to say, a countable set has countably many
          finite subsets, and a set of cardinality $\kappa$ has
          $\kappa^{\aleph_0}\geq\kappa$ countable subsets; the paper only
          needs the finite case and the fact that a countable union of
          countable sets is countable;
    \item the algebraic dimension of $\lo{\Gamma}$ is $|\Gamma|$ for infinite
          $\Gamma$, because the unit vectors $e_\gamma$ form a Hamel basis,
          and a vector space spanned by $\kappa$ vectors has dimension at most
          $\kappa$;
    \item the paper writes $\max(\operatorname{dens}(Z),\aleph_0)$ in
          theorem~3.2 only to cover finite dimensional $Z$, whose density
          character is finite.
\end{enumerate}

\begin{exercise}[*] Show that every element $x\in\ellone{\Gamma}$ has countable
    support. (Hint: for each $n$ the set $\{\gamma:|x_\gamma|>1/n\}$ is
    finite.) Deduce that every $x\in\ellone{\omega_1}$ lies in $\ellone{\delta}$
    for some countable ordinal $\delta$, using exercise~\ref{BoundedSubsets}
    below.
\end{exercise}

\subsection{The axiom of choice}\label{SecChoice}

The paper works, as most of functional analysis does, in ZFC: the
Zermelo--Fraenkel axioms with the axiom of choice. Choice enters in three
forms. The \emph{well-ordering theorem}: every set can be enumerated by an
ordinal, that is, written as $\{\gamma_\alpha:\alpha<\kappa\}$ with distinct
$\gamma_\alpha$; the paper does this with $\Gamma$ in the proof of
theorem~6.7, and with $\kappa=|\Gamma|$ one may take the enumeration to be
indexed by the cardinal itself. \emph{Zorn's lemma}: every vector space has a
Hamel basis, used throughout. \emph{Countable choice}: a countable union of
countable sets is countable, used whenever the paper says that the union of
countably many countable supports is countable, and in the proof that a
countable set of countable ordinals has a countable supremum. None of these
uses is delicate, and the reader may simply take all three as basic facts.

%-------------------------------------------------------------------------------
\section{The first uncountable ordinal}\label{SecOmegaOne}

\subsection{Basic properties of $\omega_1$}

\begin{lemma}\label{BoundedSubsets} Every countable subset of $\omega_1$ is
    bounded in $\omega_1$: if $M\subset\omega_1$ is countable, then
    $\sup M<\omega_1$.
\end{lemma}
\begin{proof} $\sup M=\bigcup M$ is a countable union of countable sets, hence
    a countable ordinal.
\end{proof}

This is the single most important property of $\omega_1$. It says that one
cannot reach $\omega_1$ in countably many steps: $\omega_1$ has
\emph{uncountable cofinality}. Contrast this with every limit ordinal
\emph{below} $\omega_1$, which can be reached in countably many steps.

\begin{lemma}[ladders]\label{Ladders} Every limit ordinal $\delta<\omega_1$ is
    the supremum of a strictly increasing sequence ${(\delta_n)}_{n\in\NN}$ of
    ordinals below $\delta$.
\end{lemma}
\begin{proof} Since $\delta$ is countable, enumerate it as $\{\beta_n:n\in\NN\}$.
    Put $\delta_0=\beta_0$ and $\delta_{n+1}=\max(\delta_n+1,\beta_{n+1})$;
    this is below $\delta$ because $\delta$ is a limit. The sequence is
    strictly increasing and $\beta_n\leq\delta_n$ for all $n$, so its
    supremum is $\delta$.
\end{proof}

Such a sequence is called a \emph{ladder} at $\delta$, and a choice of a
ladder at every limit ordinal $\delta<\omega_1$, or at every $\delta$ in a
given set $E$ of limit ordinals, is a \emph{ladder system}. The ladder spaces
of section 7 of the paper are built from a ladder system: the vector
$v_\delta=\sum_n 2^{-n}e_{\delta_n}$ ``climbs'' the ladder at $\delta$.
Two ladders at different limit ordinals $\alpha<\beta$ have only finitely
many rungs in common, because the rungs of the ladder at $\beta$ eventually
exceed $\alpha$.

\begin{exercise} Show that for every $\delta<\omega_1$ the subspace
    $\ellone{\delta}$ of $\ellone{\omega_1}$ is separable, and that
    $\ellone{\omega_1}$ is not. Show that
    $\ellone{\omega_1}=\bigcup_{\delta<\omega_1}\ellone{\delta}$, while
    $\ellone{\delta}\neq\bigcup_{\delta'<\delta}\ellone{\delta'}$ for every
    limit ordinal $\delta<\omega_1$. (Hint: lemma~\ref{Ladders}.) So the chain
    ${(\ellone{\delta})}_{\delta\leq\omega_1}$ is continuous at the top and
    discontinuous at every limit below; this is the situation the paper's
    condition $(b)$ of theorem~6.7 rules out for the pieces $P_\delta$.
\end{exercise}

\subsection{Closed unbounded sets}

\begin{definition} A set $C\subset\omega_1$ is \emph{closed} if the supremum
    of every non-empty countable subset of $C$ belongs to $C$, equivalently,
    if the supremum of every strictly increasing sequence in $C$ belongs to
    $C$. It is \emph{unbounded} if for every $\beta<\omega_1$ there is
    $\delta\in C$ with $\delta>\beta$. A \emph{closed unbounded set}, or
    \emph{club}, is a set which is both.
\end{definition}

The word ``closed'' is the topological one: in the order topology on
$\omega_1$ the closed unbounded sets are the closed sets which are unbounded.
A club is to be thought of as a \emph{large} set: it contains ``almost all''
countable ordinals in a sense made precise by the next lemmas. The simplest
club is $\omega_1$ itself. The next one is the set $\Lim$ of limit ordinals
below $\omega_1$: it is closed because a supremum of limit ordinals is a
limit ordinal (a supremum of a strictly increasing sequence is never a
successor), and unbounded because $\beta+\omega$ is a limit above $\beta$.

\begin{lemma}\label{ClosurePoints} Let $h:\omega_1\to\omega_1$ be any
    function. The set $C_h$ of ordinals $\delta<\omega_1$ such that
    $h(\alpha)<\delta$ for all $\alpha<\delta$ is a club. We say that the
    elements of $C_h$ are \emph{closed under $h$}.
\end{lemma}
\begin{proof} Closed: let $\delta$ be the supremum of a strictly increasing
    sequence ${(\delta_n)}$ in $C_h$ and $\alpha<\delta$. Then $\alpha<\delta_n$
    for some $n$, so $h(\alpha)<\delta_n\leq\delta$. Unbounded: given $\beta$,
    put $\delta_0=\beta+1$ and
    $\delta_{n+1}=\sup\{h(\alpha)+1:\alpha<\delta_n\}\cup\{\delta_n+1\}$,
    a countable ordinal by lemma~\ref{BoundedSubsets}. The sequence is strictly
    increasing; let $\delta$ be its supremum, again countable. If
    $\alpha<\delta$, then $\alpha<\delta_n$ for some $n$, so
    $h(\alpha)<\delta_{n+1}<\delta$. Thus, $\delta\in C_h$ and $\delta>\beta$.
\end{proof}

\begin{lemma}\label{IntersectionOfClubs} The intersection of two clubs is a
    club. Consequently, the intersection of finitely many clubs is a club.
\end{lemma}
\begin{proof} Let $C,D$ be clubs. Their intersection is closed, being an
    intersection of closed sets. To see that it is unbounded, let
    $\beta<\omega_1$ and choose alternately
    $c_0<d_0<c_1<d_1<\cdots$ with $c_n\in C$, $d_n\in D$ and $c_0>\beta$, which
    is possible because both sets are unbounded. The two sequences have the
    same supremum $\delta$, which is countable by lemma~\ref{BoundedSubsets},
    and $\delta\in C\cap D$ because both sets are closed.
\end{proof}

\begin{lemma}\label{ContinuousEnumeration} Let $C$ be a club and
    ${(c_\alpha)}_{\alpha<\omega_1}$ its increasing enumeration. Then
    $c_\alpha=\sup_{\alpha'<\alpha}c_{\alpha'}$ for every limit ordinal
    $\alpha<\omega_1$.
\end{lemma}
\begin{proof} The enumeration exists because $C$ is a well-ordered set; it is
    indexed by $\omega_1$ because $C$ is uncountable (unbounded) and every
    proper initial segment of $C$ is bounded, hence countable. Let $\alpha$ be
    a limit and $\gamma=\sup_{\alpha'<\alpha}c_{\alpha'}$. Then $\gamma\in C$
    because $C$ is closed, $\gamma\geq c_{\alpha'}$ for all $\alpha'<\alpha$,
    and $\gamma$ is not among the $c_{\alpha'}$ with $\alpha'<\alpha$, since
    the sequence has no largest term. Hence $\gamma$ is the least element of
    $C$ above all $c_{\alpha'}$, $\alpha'<\alpha$, which is $c_\alpha$.
\end{proof}

\begin{exercise}[*] Let $f:\omega_1\to\omega_1$ be a bijection. Show that
    $\{\delta<\omega_1:f(\delta)=\delta\}$, where $f(\delta)$ denotes the image
    of the set $\delta=\{\beta:\beta<\delta\}$, is a club. (Hint: apply
    lemma~\ref{ClosurePoints} to $f$ and to $f^{-1}$ and intersect.) This is
    the mechanism behind proposition~7.3 of the paper and behind the remark
    that the enumeration of the index set does not matter.
\end{exercise}

\begin{exercise} Show that the intersection of countably many clubs is a
    club. (Hint: for unboundedness build a sequence which visits every club
    infinitely often.) The paper needs only finite intersections.
\end{exercise}

\subsection{Stationary sets}

\begin{definition} A set $S\subset\omega_1$ is \emph{stationary} if it
    intersects every club. A set is \emph{non-stationary} if it is not
    stationary, that is, if it is disjoint from some club.
\end{definition}

If clubs are the large sets, then non-stationary sets are the small ones,
and stationary sets are the sets which are not small. Every club is
stationary by lemma~\ref{IntersectionOfClubs}; in particular $\Lim$ is
stationary. Every countable set is non-stationary, being disjoint from the
club of ordinals above its supremum. A set containing a club is stationary,
but a stationary set need not contain a club: Solovay's theorem below
provides stationary sets whose complements are stationary too, and neither
of them contains a club.

\begin{lemma}\label{NonStationaryIdeal} A subset of a non-stationary set is
    non-stationary, and the union of two non-stationary sets is
    non-stationary. If $S$ is stationary and $C$ is a club, then $S\cap C$ is
    stationary.
\end{lemma}
\begin{proof} The first claim is obvious. If $S\cap C=\varnothing$ and
    $S'\cap C'=\varnothing$ for clubs $C,C'$, then $(S\cup S')\cap(C\cap C')=\varnothing$
    and $C\cap C'$ is a club. For the last claim, if $D$ is a club, then
    $(S\cap C)\cap D=S\cap(C\cap D)\neq\varnothing$.
\end{proof}

The lemma says that the non-stationary sets form an \emph{ideal}, and
therefore the relation ``$E\triangle E'$ is non-stationary'' is an
equivalence relation on subsets of $\omega_1$: this is what ``$\Delta(P)$
modulo non-stationary sets'' means in section 7 of the paper. Two sets which
differ by a non-stationary set are stationary or non-stationary together.

\begin{theorem}[Solovay; for $\omega_1$ already Ulam]\label{Solovay} There
    is a family ${(S_\xi)}_{\xi<\omega_1}$ of pairwise disjoint stationary
    subsets of $\omega_1$ with union $\omega_1$. More generally, every
    stationary subset of $\omega_1$ is the union of $\aleph_1$ pairwise
    disjoint stationary sets.
\end{theorem}

We do not prove this; it is theorem~8.10 of Jech's book. The paper uses it
in the following way. Intersecting each $S_\xi$ with the club $\Lim$ we may
assume that all $S_\xi$ consist of limit ordinals (lemma~\ref{NonStationaryIdeal}).
For $X\subset\omega_1$ put $E_X=\bigcup_{\xi\in X}S_\xi$. If $X\neq Y$, pick
$\xi\in X\triangle Y$; then $S_\xi\subset E_X\triangle E_Y$, so
$E_X\triangle E_Y$ is stationary, and in particular $E_X\neq E_Y$ modulo
non-stationary sets. There are $2^{\aleph_1}$ sets $X$, hence $2^{\aleph_1}$
pairwise inequivalent sets $E_X$, and only $E_\varnothing=\varnothing$ is
non-stationary.

\subsection{Fodor's lemma, for orientation}

The paper does not use the following theorem, but it is the first thing one
meets after clubs and stationary sets, and it explains why ladder systems on
stationary sets are so rigid. A function $g:S\to\omega_1$ on a set $S$ of
non-zero ordinals is \emph{regressive} if $g(\alpha)<\alpha$ for all
$\alpha\in S$.

\begin{theorem}[Fodor] If $S\subset\omega_1$ is stationary and $g:S\to\omega_1$
    is regressive, then $g$ is constant on some stationary subset of $S$.
\end{theorem}

\begin{example} Let $E\subset\Lim$ be stationary and let a ladder $L_\alpha$
    at $\alpha$ be chosen for every $\alpha\in E$. Then the ladders cannot be
    pairwise disjoint: the function $g(\alpha)=\min L_\alpha$ is regressive on
    $E$, so by Fodor's lemma it takes the same value $\gamma_0$ at two
    different $\alpha$, and the corresponding ladders share the rung
    $\gamma_0$. For a non-stationary $E$ pairwise disjoint ladders can exist
    (exercise). This is the kind of phenomenon behind the fact that the
    ladder space $P_E$ of the paper is projective exactly when $E$ is
    non-stationary, although the paper proves that fact without Fodor's
    lemma.
\end{example}

\begin{exercise} Let $E$ be a non-stationary set of limit ordinals. Show that
    there are pairwise disjoint ladders $L_\alpha$, $\alpha\in E$. (Hint: let
    $C$ be a club disjoint from $E$; for $\alpha\in E$ let $c(\alpha)$ be the
    largest element of $C$ below $\alpha$, which exists because $C$ is closed
    and $\alpha\notin C$, and choose the ladder at $\alpha$ inside the
    interval $(c(\alpha),\alpha)$; distinct $\alpha$ with the same $c(\alpha)$
    lie in the same interval $(c(\alpha),c^+)$, where $c^+$ is the next element
    of $C$, and there are only countably many of them, so their ladders can
    be chosen disjoint by an inductive argument.)
\end{exercise}

%-------------------------------------------------------------------------------
\section{Reading guide to sections 6 and 7 of the paper}\label{SecReading}

\subsection*{Theorem 6.1, the transfinite back-and-forth}

The data are a filtration ${(\Gamma_\delta)}_{\delta\leq\lambda}$ with
countable blocks $N_\delta$ and the pieces $P_\delta=P\cap\ellone{\Gamma_\delta}$.
The set-theoretic content is: a transfinite recursion of length $\lambda$
defining the operators $V_\delta$, and a transfinite induction proving
$T_\delta(P_\delta)=\lo{\Gamma_\delta}$. The successor step is analysis. The
limit step is pure bookkeeping, and it is exactly where the continuity
condition $(ii)$ is used: at a limit $\mu$ the piece $P_\mu$ must be the
union of the earlier pieces, because the operators $T_\delta$ are compatible
and the images $\lo{\Gamma_\delta}$ have union $\lo{\Gamma_\mu}$. The
operator $W_\delta$ is a sum over all $\delta'<\delta$, possibly uncountably
many; it converges because every $x\in\ellone{\Gamma_\delta}$ has countable
support and meets only countably many blocks (exercise in
section~\ref{SecOrdinals}).

\subsection*{Closed pairs and closing off}

Fix the retraction $\sigma:P\to\lo{S}$, $\pi:\lo{S}\to P$ and the
approximating sequences $p^\gamma_n\to e_\gamma$. A \emph{closed pair}
$(A,B)$ is a pair of index sets which is closed under three operations:
taking the support of $\pi(e_s)$ for $s\in B$, taking the support of
$\sigma(x)$ for $x\in P_A$, and taking the supports of the $p^\gamma_n$ for
$\gamma\in A$. This is the same idea as a subgroup generated by a set, or a
Skolem hull in logic: given any sets $A_0,B_0$, apply the operations, collect
the new elements, apply them again, and so on $\omega$ times; the union is
closed under the operations, because each element of it entered at some
finite stage and its images entered at the next stage. The only subtlety,
which the paper points out, is the second operation: an element of $P_{A'}$
need not lie in any $P_{A_k}$ (its support may spread over all stages), so
closedness under $\sigma$ is proved by continuity of $\sigma$ and density of
$\bigcup_kP_{A_k}$ rather than by a support argument.

Lemma~6.5 of the paper says that if one starts from a closed pair and adds a
countable set, then $\omega$ steps of closing off add only countably many
points. The reason is lemma~6.4: over a closed pair, a countable extension
of $A$ changes $P_A$ by a quotient of countable dimension, so only countably
many new supports are needed at each stage. This is the analogue of the
downward L\"owenheim--Skolem theorem: a countable set generates a countable
closed structure.

\subsection*{Theorem 6.7, the chain of closed pairs}

Enumerate $\Gamma=\{\gamma_\alpha:\alpha<\kappa\}$ (well-ordering theorem)
and define closed pairs by transfinite recursion: add $\gamma_\alpha$ and
close off at successor steps, take unions at limits. Lemma~6.6 says that a
union of closed pairs along a chain is closed, and this is what makes the
limit step work and gives the continuity condition for free. The resulting
family ${(A_\alpha)}_{\alpha\leq\kappa}$ is a filtration with countable blocks,
and the three hypotheses of theorem~6.1 are verified by the three lemmas on
closed pairs. Nothing about the size of $\Gamma$ matters here; this is why
the paper no longer needs an induction on cardinals.

\subsection*{Theorem 7.2, the criterion for $\omega_1$}

Here $\Gamma=\omega_1$ and $P_\delta=P\cap\ellone{\delta}$. The theorem says:
$P$ is extremely projective if and only if the pieces have countable
dimension and the chain ${(P_\delta)}$ is continuous on a club, that is, the
set $\Delta(P)$ of limit ordinals at which it jumps is non-stationary.

For the forward direction one runs the chain of closed pairs with
$\gamma_\alpha=\alpha$. Each $A_\alpha$ is a countable subset of $\omega_1$
(countable increments, countable unions of countable sets). Define
$f(\alpha)=\sup A_\alpha+1$ and $g(\beta)=$ the first $\alpha$ with
$\beta\in A_\alpha$. On the club of limit ordinals closed under $f$ and $g$
(lemmas~\ref{ClosurePoints} and~\ref{IntersectionOfClubs}) one has
$A_\delta=\delta$: closure under $f$ gives $A_\delta\subset\delta$, closure
under $g$ gives $\delta\subset A_\delta$. At such $\delta$ the union lemma
turns $P_\delta$ into a union of earlier pieces, so $\delta\notin\Delta(P)$.

For the converse one needs a filtration. Let
\[
h(\gamma)=\sup\textstyle\bigcup_n\supp(p^\gamma_n)+1 .
\]
The club $C_1$ of ordinals closed under $h$ consists of ordinals at
which the piece is dense; a club $C_2$ misses $\Delta(P)$; the enumeration
of $(C_1\cap C_2)\cup\{0\}$ is continuous
(lemma~\ref{ContinuousEnumeration}) and
defines a filtration of $\omega_1$ with countable blocks. Continuity of the
chain of pieces along it holds because its limit points are limit ordinals
in $C_2$, and the quotients are countable because the pieces are. Theorem~6.7
finishes the proof.

\subsection*{Proposition 7.3, invariance}

Given a bounded bijection $T$ of $\ellone{\omega_1}$, the function
$F(\delta)=\sup_{\beta<\delta}(\sup\supp T(e_\beta)+1)$ records how far $T$
moves the coordinates below $\delta$. On the club of limit ordinals closed
under $F$ and under the analogous $G$ for $T^{-1}$, $T$ maps $\ellone{\delta}$
onto itself, hence maps the chain of pieces of $P$ onto that of $T(P)$, and
jumps correspond. So $\Delta(P)$ and $\Delta(T(P))$ agree on a club, that is,
modulo a non-stationary set. Applied to a permutation of the index set this
shows that the choice of enumeration is immaterial.

\subsection*{Theorem 7.4, ladder spaces}

For a set $E$ of limit ordinals the ladder space is
$P_E=\lo{\omega_1}+\Lin\{v_\alpha:\alpha\in E\}$ with $v_\alpha$ climbing a
ladder at $\alpha$ (lemma~\ref{Ladders}). The computation
$\Delta(P_E)=E$ uses only that two ladders share finitely many rungs. Then
theorem~7.2 says that $P_E$ is extremely projective exactly when $E$ is
non-stationary, proposition~7.3 says that $P_E\not\cong P_{E'}$ when
$E\triangle E'$ is stationary, and Solovay's theorem~\ref{Solovay} produces
$2^{\aleph_1}$ sets $E_X$ with pairwise stationary symmetric differences.

\subsection*{Remark 7.5, almost free abelian groups}

An abelian group is free if it has a basis; it is $\aleph_1$-free if all its
countable subgroups are free. Pontryagin's criterion and Hill's lemma say
that a group of cardinality $\aleph_1$ is free if and only if it is the union
of a continuous $\omega_1$-chain of countable free subgroups with free
successive quotients; Eklof's invariant is the set of limit stages at which
such a chain cannot be continued with a free quotient, taken modulo
non-stationary sets; and ladder systems on stationary sets give
$\aleph_1$-free groups which are not free. The paper's theorem~6.7,
definition~7.1 and theorem~7.4 are the exact counterparts, with ``countable
dimension'' in place of ``free''. The reader who wants to see where this
leads above $\aleph_1$ may consult chapters IV and VII of Eklof and Mekler.

%-------------------------------------------------------------------------------
\section{Glossary}\label{SecGlossary}

\begin{description}[style=nextline,leftmargin=1em]
    \item[well-order] a linear order in which every non-empty subset has a
          least element (section~\ref{SecOrdinals}).
    \item[ordinal] a canonical well-order; the set of all smaller ordinals
          (section~\ref{SecOrdinals}).
    \item[successor ordinal, limit ordinal] $\alpha+1$, respectively a non-zero
          ordinal which is not a successor; a limit is the union of the
          ordinals below it (section~\ref{SecOrdinals}).
    \item[$\omega$] the first infinite ordinal, the set of natural numbers;
          as a cardinal, $\aleph_0$.
    \item[$\omega_1$] the first uncountable ordinal, the set of countable
          ordinals; as a cardinal, $\aleph_1$ (section~\ref{SecCardinals}).
    \item[$\omega_2$, $\aleph_2$] the next initial ordinal and cardinal;
          mentioned in the paper only in the open questions.
    \item[transfinite induction, recursion] proving a property, respectively
          defining an object, at $\alpha$ from the ordinals below $\alpha$
          (section~\ref{SecOrdinals}).
    \item[enumeration by an ordinal] a bijection $\alpha\mapsto\gamma_\alpha$
          from an ordinal $\kappa$ onto a set $\Gamma$; exists by the
          well-ordering theorem (section~\ref{SecChoice}).
    \item[filtration, block] an increasing continuous family of subsets of
          $\Gamma$ starting at $\varnothing$ and ending at $\Gamma$; the
          differences of consecutive members (section~\ref{SecOrdinals}).
    \item[cardinality of $\Gamma$, $|\Gamma|$] the least ordinal in bijection
          with $\Gamma$ (section~\ref{SecCardinals}).
    \item[density character] the least cardinality of a dense subset; for
          $\ellone{\Gamma}$ with $\Gamma$ infinite it is $|\Gamma|$.
    \item[closed unbounded set, club] a subset of $\omega_1$ closed under
          countable suprema and unbounded (section~\ref{SecOmegaOne}).
    \item[closed under $h$] an ordinal $\delta$ with $h(\alpha)<\delta$ for all
          $\alpha<\delta$; these ordinals form a club (lemma~\ref{ClosurePoints}).
    \item[stationary, non-stationary] meeting every club, respectively
          disjoint from some club (section~\ref{SecOmegaOne}).
    \item[modulo non-stationary sets] the equivalence relation
          ``$E\triangle E'$ is non-stationary'' (lemma~\ref{NonStationaryIdeal}).
    \item[ladder, ladder system] a strictly increasing $\omega$-sequence
          converging to a limit ordinal; a choice of one at every limit
          ordinal of a given set (lemma~\ref{Ladders}).
    \item[Solovay's theorem] $\omega_1$ splits into $\aleph_1$ disjoint
          stationary sets (theorem~\ref{Solovay}).
    \item[Fodor's lemma] a regressive function on a stationary set is
          constant on a stationary set; not used in the paper
          (section~\ref{SecOmegaOne}).
    \item[continuum hypothesis] $2^{\aleph_0}=\aleph_1$; not used in the paper.
\end{description}

\begin{thebibliography}{9}
    \bibitem{Jech} \textit{T. Jech} Set theory. The third millennium edition,
    revised and expanded, Springer, Berlin (2003). Chapters 2 (ordinals),
    3 (cardinals), 5 (the axiom of choice), 8 (stationary sets).
    \bibitem{HrbacekJech} \textit{K. Hrbacek, T. Jech} Introduction to set
    theory, third edition, Marcel Dekker, New York (1999).
    \bibitem{Kunen} \textit{K. Kunen} Set theory, College Publications,
    London (2011).
    \bibitem{EklofMekler} \textit{P. C. Eklof, A. H. Mekler} Almost free
    modules. Set-theoretic methods, revised edition, North-Holland, Amsterdam
    (2002).
\end{thebibliography}

\end{document}
